\scanpage{101}
que $\underline{\mathrm{Pic}}^\tau_{X/k}$ est propre, on est ramené à le prouver pour une \uline{courbe} normale, où c'est connu. Alors un sous-grp. unipotent est néc. fini, ce qui achève de prouver l'assertion \ldots

\medskip\centerline{\rule{1.2cm}{0.4pt}}

\uline{Démonstration directe et indépendante de la \unsure{car.|con.} du fait que si $X$ est} \inserted{\uline{géom$^t$}} \uline{normal et projectif} $/k$, \deleted{\unsure{le grp.}}\ $\underline{\mathrm{Pic}}^0_{X/k}$ est propre.

On utilise le théorème de structure des groupes algébriques commutatifs, on est ramené à prouver que \unsure{tout|toute} morphisme $\operatorname{Spec}(k[t,t^{-1}])\to\underline{\mathrm{Pic}}^0_{X/k}$ est constant, \deleted{et il suffit [$k$ clos alg. bien \inserted{sûr.}] de prouver que l'application \unsure{sous-jacente|sans augment} est trivial. \unsure{Cela|Ceci} provient de ceci}

\uline{Lemme} \ Supposons $X$ préschéma loc noeth - normal \inserted{existe un fermé $Y$ de codim $\geq2$ \deleted{\unsure{dont|de}} tel que $U:X-Y$ soit régulier}
\margin{Certainement inutile, en procédant de façon convenable}
Alors \deleted{pour tout}
\[
\mathrm{Pic}(X)\longrightarrow\mathrm{Pic}\,X[t,t^{-1}]
\]
est \deleted{bijectif} \uline{surjectif}.

\deleted{C'est injectif, car $X[t]$ a une section sur $X$ (section nulle).} Soit donc $\underline L$ faisceau inversible sur $X'=X[t,t^{-1}]$, \deleted{soit $U$ l'ens. des pts réguliers de $X$} \inserted{on \unsure{sait} que} \inserted{posant $U'=X'(U)$} $\underline L\mid U'\deleted{\unsure{[t]}}$ provient d'un faisceau inversible $\underline L_0$ sur $U$, qui est défini à l'aide d'un diviseur $D_0$ sur $U$.

Son adhérence $D_1=\overline D_0$ est tel que son image inverse sur $X'\deleted{\unsure{[t]}}$ \deleted{\unsure{est}} soit loc. principale. Pour raison de platitude, il en est de $=$ de $D_1$, d'où facilement la conclusion.
% archive sheet number: 100

